\scanpage{263}
\uline{Ex} \quad $M=S[-i]$, on \unsure{trouve}
\[
\begin{aligned}
\unsure{\kappa|k|H}_i(R^*\unsure{\varphi|\rho}_N(\tau_N(P))&\simeq\mathbb R\operatorname{Hom}(\tau_N(S[-i]),P)\\
&\simeq\mathbb R\operatorname{Hom}((\tau_{N-i}S)[-i],P)\\
&\simeq R\operatorname{Hom}^i(\tau_{N-i}S,P).
\end{aligned}
\]
donc
\[
(R^n\varphi_N(\tau_N P))_i=\mathbb E\mathrm{xt}^n((\tau_{N-i}S)[-i],P)
\]
\uline{NB} C'est valable aussi en remplaçant $\varphi_N$ par $\varphi_{N',N}$ $(N'\leq N)$ pourvu que $i\geq N'$ (si $i<N'$, on trouve 0 bien sûr !)

\deleted{\uline{NB} Supposons $P\in\operatorname{Im\,ess}\varphi_{N',N}$ \quad $(N'\leq N)$

i.e. $\deleted{\unsure{\varphi_{N'}}}P\xrightarrow{\sim}\varphi_{N',N}\tau_N P=R^0\varphi_{N',N}(\tau_N P)$

Alors $R^n(\varphi_{N',N}(\tau_N P))=0$ si $n>0$

i.e. $\operatorname{Ext}^n(\tau_{N-i}(S)[-i],P)=0$ si $\substack{i\geq N'\\ n>0}$

En effet, cela revient à}

\deleted{\[
R^{\unsure{n|*}}\unsure{\varphi_N}\varphi_N\unsure{\tau(P)}_{\substack{\rotatebox{90}{$\in$}\\\mathcal M_N}}
\]}

Soit $C$ une catégorie abélienne, \inserted{$\unsure{N'_0|N_0^*|N_0}\leq N$ ds $\mathbb Z$} et
\[
C\xrightarrow{i}\mathcal M_{\unsure{N_0|N'_0}}
\]
un foncteur \uline{exact}, tel que
\[
i(C)\subset\operatorname{Im\,ess}\varphi_{N'_0,N}
\]
i.e
\[
i(\unsure{\deleted{M}|\underline M|M})\xrightarrow{\sim}\varphi_{N'_0,N}\tau_{N,N'_0}i(M)
\]
Je dis qu'alors
\[
(R^n\varphi_{N'_0,N})(\tau_{N,N'_0}\circ i)=0\quad\text{si }n>0
\]
En effet, comme $i\simeq\varphi_{N'_0,N}\tau_{N,N'_0}i$ on en déduit
\diagram{\begin{tikzcd}[ampersand replacement=\&,column sep=small]
R^\alpha i\arrow[d,phantom,"\rotatebox{90}{$\simeq$}"] \& {}\simeq R^\alpha(\varphi_{N'_0,N}\tau_{N,N'_0})\arrow[d,phantom,"\rotatebox{90}{$\simeq$}"]\\
i\quad\text{car }i\text{ exact}\& (R^*\varphi_{N'_0,N})(\tau_{N,N'_0}\circ i)\quad\text{car }\tau_{N,N'_0}\text{ est exact}
\end{tikzcd}}
% archive sheet number: 262
